% =============================================================================
%  CHAPTER 10 -- MULTIPLE TESTING AND DATA SNOOPING
% =============================================================================
\section{Multiple testing, data snooping, and what survives}
\label{ch:multipletesting}

This is the chapter that decides whether the project's results mean anything. It is also
the one most often skipped, which is why most published backtests are wrong.

\subsection{The arithmetic of searching}
\label{sec:mt:arithmetic}

Suppose you test $M$ strategies that are all pure noise, each with an independent test at
level $\alpha$. The probability that \emph{at least one} shows a ``significant'' result is
\begin{equation}
  P(\text{at least one false discovery}) \;=\; 1-(1-\alpha)^{M}
  \;\overset{\alpha M\ll1}{\approx}\; \alpha M .
  \label{eq:familywise}
\end{equation}
For $\alpha=0.05$: $M=10\Rightarrow36\%$, $M=44\Rightarrow89\%$, $M=100\Rightarrow99.4\%$.
The number audited in Phase~\ref{ch:phase10} is $M=44$ (the product of 4 baskets
$\times$ 11 hedge configurations actually traded and reported across Phases 4--9). At
$M=44$ a single 5\%-level ``discovery'' is worth essentially nothing.

\paragraph{The expected maximum is the sharper statement.} If $t_1,\dots,t_M$ are i.i.d.\
standard normal $t$-statistics under the null, the distribution of the maximum satisfies,
for large $M$,
\begin{equation}
  \E\bigl[\max_{m\le M} t_m\bigr] \;\approx\; \sqrt{2\ln M},
  \qquad
  \Var\bigl[\max_m t_m\bigr] \;\approx\; \frac{\pi^2}{6\ln M},
  \label{eq:maxnormal}
\end{equation}
(the Gumbel extreme-value limit; $\sqrt{2\ln M}$ comes from inverting the normal tail
$1-\Phi(x)\approx\phi(x)/x$ at $x\approx\sqrt{2\ln M}$ where $M(1-\Phi(x))\approx1$).
For $M=44$: $\E[\max t]\approx\sqrt{2\ln 44}=2.75$.

\textbf{Now compare with what we measured.} The best of our 44 strategies has a HAC
$t$-statistic for ``mean daily net P\&L $>0$'' of $\mathbf{1.57}$
(Table~\ref{tab:multtest_audit}: \code{GS::OLS-w180}, net Sharpe $0.346$, mean daily net
$1.03$\,bps, raw one-sided $p=0.059$). The expected maximum of \emph{44 pure-noise}
strategies is $2.75$. Our best is well \emph{below} the noise benchmark.

Two honest qualifications, both measured rather than assumed
(Table~\ref{tab:power_dsr}). (i) The 44 strategies are correlated --- they share baskets,
windows, and the same cost model --- so the \emph{effective} number of independent trials
is smaller than 44. The mean pairwise correlation of the daily net P\&L series is only
$0.128$, and the participation ratio of the correlation spectrum,
$M_{\text{eff}}=(\sum_i\lambda_i)^2/\sum_i\lambda_i^2$, gives $M_{\text{eff}}=10.7$ ---
which still implies $\E[\max t]\approx\sqrt{2\ln 10.7}=2.18$, comfortably above the
observed $1.57$. (ii) The strategies are not exactly null: they trade real data with a
real (if small) signal. But the direction of the comparison is unambiguous and does not
depend on any correction methodology: \emph{the best of the searched set is not impressive
relative to what searching produces}.

\paragraph{The deflated Sharpe ratio says the same thing.} Applying
\eqref{eq:sr0}--\eqref{eq:dsr} with $N=M=44$, cross-trial Sharpe variance estimated from
the 44 realized Sharpers, and the best strategy's own daily P\&L moments (skewness
$+0.473$, excess kurtosis $+13.4$ --- the fat tails of a spread book that mostly earns
small amounts and occasionally loses large ones), the expected best annualised Sharpe
under the null is $\mathrm{SR}_0 = 0.474$, and
\begin{equation}
  \mathrm{DSR} \;=\; 0.305 ,
  \label{eq:dsr:measured}
\end{equation}
i.e.\ a 30.5\% probability that the true Sharpe of the best-of-44 exceeds what selection
alone would produce --- against the conventional $>0.95$ bar. The stationary-bootstrap
95\% percentile interval for that Sharpe is $[-0.058,\ +0.753]$, which contains zero.
Note that the fat tails \emph{widen} the denominator of \eqref{eq:dsr} through the
$(\hat\gamma_4-1)/4$ term, so non-normality makes the case worse, not better --- the
opposite of the direction in which a kurtotic return series is usually (wrongly) said to
help a Sharpe ratio.

\subsection{Bonferroni and its refinements}
\label{sec:mt:bonferroni}

Boole's inequality, $P(\bigcup_m A_m)\le\sum_m P(A_m)$, immediately gives the
\emph{family-wise error rate} (FWER) control: testing each of $M$ hypotheses at level
$\alpha/M$ guarantees $P(\text{any false rejection})\le\alpha$. Equivalently, compare the
raw $p$-value against $\alpha/M$:
\begin{equation}
  \alpha_{\text{Bonf}} \;=\; \frac{0.05}{44} \;=\; 0.001136,
  \qquad
  p^{\text{Bonf}}_m \;=\; \min\bigl(1,\ M\cdot p_m\bigr).
  \label{eq:bonferroni}
\end{equation}
Applied to the best strategy: $p^{\text{Bonf}}=\min(1,44\times0.0587)=1.0$. Number of
strategies surviving: \textbf{0}.

Bonferroni is conservative when tests are positively correlated (which ours are), because
Boole's inequality is tight only for disjoint events. Refinements, in increasing order of
sophistication:
\begin{itemize}
  \item \textbf{Holm--Bonferroni} (step-down): order the $p$-values
    $p_{(1)}\le\dots\le p_{(M)}$ and compare $p_{(j)}$ against $\alpha/(M-j+1)$, stopping
    at the first non-rejection. Strictly more powerful than Bonferroni at the same FWER,
    with no assumption about dependence. Applying it here changes nothing: the smallest
    $p$ is $0.0587$, and the Holm bar for $j=1$ is $0.05/44=0.00114$.
  \item \textbf{\v{S}idák}: $\alpha_{\text{\v S}}=1-(1-\alpha)^{1/M}$, exact under
    independence, slightly less conservative than Bonferroni. Here
    $1-0.95^{1/44}=0.00116$ --- no practical difference.
  \item \textbf{Benjamini--Hochberg} (false \emph{discovery} rate): controls
    $\E[\text{false rejections}/\text{rejections}]$ rather than the probability of any.
    More powerful, appropriate when you expect many true effects and want a ranked list.
    Here there are no rejections at all, so FDR control is moot --- but it is the right
    tool for a 500-basket universe scan, which is why it is listed as future work.
\end{itemize}

\subsection{Bootstrap inference for dependent data}
\label{sec:mt:bootstrap}

White's Reality Check and Hansen's SPA both need the null distribution of a maximum over
$M$ correlated series, which has no closed form. It is obtained by resampling --- but an
i.i.d.\ bootstrap is invalid here because the daily P\&L series are autocorrelated
(the mean reversion itself) and cross-correlated (shared baskets). Resampling individual
days destroys the dependence and produces over-confident intervals.

\paragraph{Block bootstrap.} Resample contiguous blocks of length $L$ to preserve local
dependence. The bias: block boundaries are artificial, and the resampled series is not
stationary.

\paragraph{Stationary bootstrap} \cite{politisromano1994}, implemented in
\code{scripts/statkit/bootstrap.py}. Choose a starting index uniformly, then at each step
either continue to the next index (probability $1-p$) or jump to a new uniformly chosen
index (probability $p$), wrapping around the end of the sample:
\begin{equation}
  i_{t+1} = \begin{cases} (i_t+1)\bmod T & \text{with probability } 1-p,\\
  \ \ \text{Uniform}\{0,\dots,T-1\} & \text{with probability } p.
  \end{cases}
  \label{eq:stationarybootstrap}
\end{equation}
Block lengths are therefore \emph{geometric} with mean $1/p$, and --- this is the point ---
the resampled series is \emph{strictly stationary} (its distribution is invariant to the
starting point), unlike the fixed-block version. We use $p=1/10$, i.e.\ mean block length
10 bars, chosen to exceed the $\approx9.4$-bar half-life of the traded spread so that a
typical reversion episode survives intact inside a block. $B=1{,}200$ draws.

\paragraph{Percentile confidence intervals.} Given bootstrap replicates
$\hat\theta^{(1)},\dots,\hat\theta^{(B)}$ of a statistic, the percentile CI is
$[\hat\theta^{*(\alpha/2)},\hat\theta^{*(1-\alpha/2)}]$ --- the empirical quantiles. It is
first-order correct and requires no distributional assumption; it can undercover for
biased statistics, in which case the \emph{basic}/pivotal interval
$[2\hat\theta - \hat\theta^{*(1-\alpha/2)},\ 2\hat\theta-\hat\theta^{*(\alpha/2)}]$ is
preferred. We report percentile intervals on the Sharpe ratio.

\subsection{White's Reality Check}
\label{sec:mt:whiterc}

\cite{white2000realitycheck} frames the question correctly: not ``is strategy $m$ better
than the benchmark'' for each $m$, but
\begin{equation}
  H_0:\ \max_{m\le M}\ \E[f_{m}] \;\le\; 0
  \qquad\text{vs.}\qquad
  H_1:\ \max_m \E[f_m] > 0,
  \label{eq:rc:null}
\end{equation}
where $f_{m,t}$ is strategy $m$'s performance differential versus the benchmark (here,
zero: the net daily P\&L itself). The statistic is the best observed mean,
\begin{equation}
  V_T \;=\; \max_{m\le M}\ \bar f_m,
  \qquad \bar f_m = \frac1T\sum_{t=1}^T f_{m,t},
  \label{eq:rc:stat}
\end{equation}
and the $p$-value is the bootstrap probability, under a \emph{recentered} null, of
exceeding $V_T$:
\begin{equation}
  p \;=\; \frac{1}{B}\sum_{b=1}^{B}\mathbf 1\Bigl\{\max_m \bar f_m^{*(b)} \ \ge\ V_T\Bigr\},
  \qquad
  f_{m,t}^{*(b)} = \bigl(f_{m,t} - \bar f_m\bigr)^{*(b)} .
  \label{eq:rc:pvalue}
\end{equation}
The recentering $f_{m,t}-\bar f_m$ is what imposes the null: it forces every strategy's
bootstrap mean to be exactly zero while preserving the entire covariance structure across
strategies and lags. Because the \emph{maximum} over all $M$ is recomputed in each
bootstrap draw, the search itself is baked into the null distribution --- which is exactly
what a per-strategy test cannot do.

A technical caveat that matters for interpretation: if some strategies have
$\E[f_m]<0$ (many of ours do, badly), recentering to zero \emph{raises} them and makes the
null distribution of the max wider, which makes the test conservative. Hansen's SPA fixes
this.

\paragraph{Measured: $p_{\text{RC}} = 0.868$.} Not remotely significant. The bootstrap
distribution of the best-of-44 mean P\&L under the null exceeds our observed best in 87\%
of draws.

\subsection{Hansen's Superior Predictive Ability test}
\label{sec:mt:hansen}

\cite{hansen2005test} improves the Reality Check in three ways.
\begin{enumerate}
  \item \textbf{Studentisation.} Divide each strategy's mean by its own standard error,
    \begin{equation}
      T_{\text{SPA}} \;=\; \max_{m\le M}\ \frac{\sqrt{T}\,\bar f_m}{\hat\omega_m},
      \qquad \hat\omega_m^2 = \widehat{\mathrm{LRV}}(f_{m,t}),
      \label{eq:spa}
    \end{equation}
    so that high-variance strategies do not dominate the max simply by being noisy. This
    is a scale-invariance the Reality Check lacks.
  \item \textbf{A consistent recentering rule.} Recenter only by strategies that are
    \emph{not} significantly worse than the benchmark (an indicator $\mathbf 1\{\sqrt T
    \bar f_m > -\sqrt{2\log\log T}\,\hat\omega_m\}$), which makes the test consistent
    against alternatives while retaining correct size.
  \item \textbf{Two variants} (``lower'' and ``upper'' bounds on the $p$-value) depending
    on how the poorly performing strategies are treated.
\end{enumerate}
The SPA null is ``no strategy is superior to the benchmark''; rejecting means at least one
is. Measured: $p_{\text{SPA}}=0.638$ --- again no rejection. Hansen's test is generally
more powerful than White's when many irrelevant strategies are included, so the fact that
both agree strengthens the conclusion.

\subsection{Diebold--Mariano: comparing two forecasts}
\label{sec:mt:dm}

\cite{dieboldmariano1995}. Given two forecast-error series $e^{(1)}_t,e^{(2)}_t$ and a
loss $L(\cdot)$ (squared error for forecast accuracy, or \emph{negative P\&L} for traded
strategies), define the loss differential $d_t = L(e^{(1)}_t)-L(e^{(2)}_t)$. Under
$H_0:\E[d_t]=0$,
\begin{equation}
  \mathrm{DM} \;=\; \frac{\bar d}{\sqrt{\widehat{\mathrm{LRV}}(\bar d)/T}}
  \;\xrightarrow{\ d\ }\;\mathcal N(0,1),
  \qquad
  \widehat{\mathrm{LRV}}(\bar d) = \hat\gamma_0 + 2\sum_{\ell=1}^{L}w_\ell\hat\gamma_\ell ,
  \label{eq:dm}
\end{equation}
with $\hat\gamma_\ell = \frac1T\sum_t (d_t-\bar d)(d_{t-\ell}-\bar d)$ and Bartlett
weights $w_\ell=1-\ell/(L+1)$ --- the same HAC estimator as \eqref{eq:hac}, applied to the
differential series. The HAC step is not optional: with multi-day holds and overlapping
targets, $d_t$ is strongly autocorrelated, and using $\hat\gamma_0$ alone understates the
variance by a factor of $1+2\sum_\ell\rho(\ell)$, which can easily be $5$--$10\times$ and
turns a null difference into a ``significant'' one.

\paragraph{Small-sample correction.} \cite{harvey1997evaluating} show the DM statistic
over-rejects in small samples and propose multiplying by
$\sqrt{(T+1-2h+h(h-1)/T)/T}$ for $h$-step-ahead forecasts; with $h=1$ this is
$\approx\sqrt{(T)/T}=1$ and no correction is needed, which is our case (daily P\&L
differentials, one-step).

\paragraph{Measured: nothing is distinguishable.}
Table~\ref{tab:multtest_dm} and Figure~\ref{fig:multtest_forest} run the DM test between
the best ADBE hedge (Kalman $q=10^{-5}$) and every other ADBE variant. Mean daily
differences range from $0.03$ to $0.86$ bps with HAC standard errors of $0.09$--$1.16$
bps, giving $|\mathrm{DM}|$ between $0.04$ and $1.23$ and $p$-values between $0.22$ and
$0.96$. \emph{Every} 95\% confidence interval spans zero. So although the smooth Kalman
hedge beats the churny 120-bar OLS by $0.48$ bps/day --- an annualised difference of
$\approx12$\% on \$1 gross, which is economically enormous --- sixteen years of one basket
cannot establish it statistically. That is the power limitation of
equation~\eqref{eq:sharpetstat}, and it is the most important sentence in this chapter.

\subsection{The deflated Sharpe ratio}
\label{sec:mt:dsr}

\cite{bailey2014deflated} combine the Sharpe-ratio sampling distribution
\eqref{eq:sharpevar2} with the extreme-value expectation \eqref{eq:maxnormal} into a
single probability. Define the \emph{expected maximum Sharpe under the null} over $N$
independent trials with cross-trial Sharpe variance $V[\widehat{\mathrm{SR}}]$:
\begin{equation}
  \mathrm{SR}_0 \;=\; \sqrt{V[\widehat{\mathrm{SR}}]}\,\Bigl[(1-\gamma)\,
  \Phi^{-1}\Bigl(1-\tfrac1N\Bigr) + \gamma\,\Phi^{-1}\Bigl(1-\tfrac{1}{N e}\Bigr)\Bigr],
  \qquad \gamma = 0.5772\ \text{(Euler--Mascheroni)},
  \label{eq:sr0}
\end{equation}
and the \emph{deflated Sharpe ratio}, the probability that the true Sharpe exceeds
$\mathrm{SR}_0$ given the observed estimate, its sampling error, and the non-normality of
returns:
\begin{equation}
  \mathrm{DSR} \;=\; \Phi\!\left(
  \frac{\bigl(\widehat{\mathrm{SR}}-\mathrm{SR}_0\bigr)\sqrt{T-1}}
       {\sqrt{1-\hat\gamma_3\,\widehat{\mathrm{SR}} + \frac{\hat\gamma_4-1}{4}\,
        \widehat{\mathrm{SR}}^{\,2}}}\right),
  \label{eq:dsr}
\end{equation}
with $\hat\gamma_3$ the sample skewness, $\hat\gamma_4$ the sample (non-excess) kurtosis,
and $\widehat{\mathrm{SR}}$ in per-period units. DSR $>0.95$ is the conventional bar for
``this Sharpe is not a selection artifact''. Table~\ref{tab:power_dsr} reports the
computed values for the best strategy in the audited set, with $N=M=44$ and
$V[\widehat{\mathrm{SR}}]$ estimated as the variance of the 44 realized Sharpers.

\subsection{Power: how much data would it take?}
\label{sec:mt:power}

Inverting \eqref{eq:sharpetstat} gives the sample size needed to detect a true daily
Sharpe $\mathrm{SR}_d$ at level $\alpha$ with the critical value $z_{1-\alpha}$:
\begin{equation}
  T^{\ast} \;=\; \Bigl(\frac{z_{1-\alpha}}{\mathrm{SR}_d}\Bigr)^{2}
  \;=\; \Bigl(\frac{z_{1-\alpha}\sqrt{252}}{\mathrm{SR}_a}\Bigr)^{2}\ \text{days}.
  \label{eq:requiredT}
\end{equation}
This is the calculation that turns a negative result into a research programme. Applied to
the measured numbers (Table~\ref{tab:power_dsr}: best-of-44 annualised net Sharpe
$0.346$, daily $0.0218$; best \code{ADBE} variant $\mathrm{SR}_a=0.086$), with
$z=1.645$ for a one-sided 5\% test and $z=3.052=\Phi^{-1}(1-0.05/44)$ for the
Bonferroni-corrected level:

\begin{itemize}
  \item $\mathrm{SR}_a=0.346$ (the best of all 44, \code{GS::OLS-w180}):
    $T^\ast=(1.645/0.0218)^2=5{,}695$ days $=\mathbf{22.6}$ \textbf{years} uncorrected;
    $T^\ast=(3.052/0.0218)^2=19{,}608$ days $=\mathbf{77.8}$ \textbf{years} after
    Bonferroni. We have 15.8.
  \item $\mathrm{SR}_a=0.086$ (the best \code{ADBE} variant, i.e.\ the basket the bulk of
    this report is about): $T^\ast=\mathbf{366}$ \textbf{years} uncorrected and
    $\mathbf{1{,}259}$ \textbf{years} Bonferroni-corrected.
\end{itemize}

\begin{remark}[Honesty: the verdict, stated plainly]
No correction methodology is needed to reach the project's conclusion; the power
calculation alone reaches it. A single-basket, daily-frequency market-neutral spread
strategy with a net Sharpe below $0.1$ is \emph{unfalsifiable} on any amount of historical
data that exists: you would need 366 years of daily observations for an uncorrected
one-sided test, and 1{,}259 years once the search over 44 configurations is paid for.
This is not a
defect of our analysis, it is a property of the strategy class. The implications are
concrete and they drive Chapter~\ref{ch:conclusion}: (i) get \emph{breadth} rather than
precision --- 500 independent baskets at $\mathrm{SR}=0.07$ each is a portfolio with
$\mathrm{SR}=0.07\sqrt{500}=1.6$ by \eqref{eq:fundamentallaw}, testable in 5 years;
(ii) increase frequency, where $T$ per year is 250$\times$ larger and half-lives are
minutes rather than days; (iii) accept a lower cost base, since
\eqref{eq:costdrag} shows cost consumes 87\% of the gross edge at the current turnover.
\end{remark}

\subsection{What would change the conclusion}
\label{sec:mt:whatwould}

The audit is only as good as its hypothesis count. Three things would legitimately change
the verdict, and none of them is ``try more configurations'':
\begin{enumerate}
  \item \textbf{A pre-registered single hypothesis.} Fix the basket-construction rule,
    the estimator, the window, the thresholds and the cost model \emph{before} looking at
    any P\&L, then run once. $M=1$, so $\alpha=0.05$ and the bar is
    $T^\ast\approx5{,}700$ days for $\mathrm{SR}=0.35$ --- attainable with 23 years of
    data, or with 23 baskets over one year.
  \item \textbf{Genuinely independent breadth.} $M$ grows, but so does $T$ in the
    portfolio sense: 200 disjoint baskets give $\sqrt{200}=14\times$ the effective sample
    for the \emph{book}, and the relevant test becomes one test on the book, not 200 tests
    on its parts.
  \item \textbf{Out-of-sample in time.} The only correction that cannot be gamed is fresh
    data. Run the frozen configuration forward for two years and test once.
\end{enumerate}
Everything else --- more windows, more baskets searched post hoc, more hyper-parameter
grids --- increases $M$ faster than it increases power.
